\section{Related Literature and Contribution}
\label{sec:lit}

This paper is positioned at the intersection of three literatures that are central to regional science: work on local interaction and regional coordination, work on proximity, lock-in, and system renewal, and work on the welfare evaluation of place-based policy. The paper's aim is not to offer a general theory of regional development. It is to provide a tractable policy framework for a specific regional question: when a publicly supported hub in a thin regional ecosystem should be evaluated as a stand-alone coordination device, and when it should instead be evaluated together with a policy that builds cross-regional pipelines.

\subsection{Local Interaction, Thin Regional Ecosystems, and Coordination Devices}

A first relevant literature explains why spatial concentration and local interaction can raise productivity. In urban and regional economics, agglomeration economies are commonly rationalized through matching, learning, and sharing mechanisms \citep{DurantonPuga2004,RosenthalStrange2004}. A hub policy fits naturally into this tradition. It can be interpreted as a public coordination device that reduces local interaction frictions and raises the rate at which potentially valuable local matches are realized within a region.

This logic is reinforced by work emphasizing the distinctive role of face-to-face contact, tacit knowledge exchange, and proximate social interaction \citep{StorperVenables2004,JaffeTrajtenbergHenderson1993,RocheOettlCatalini2024}. At the same time, recent work on co-working spaces, enterprise hubs, and related collaborative venues in rural and peripheral settings suggests that these devices are highly context-dependent \citep{VoglAkhavan2022,MerrellRoweCowieGkartzios2021,MerrellPhillipsonGortonCowie2022,OECD2025}. Thin and nonmetropolitan regions may benefit from organized interaction, but the gains are neither automatic nor uniform.

The present paper takes this literature as a starting point, but shifts the question. Rather than asking whether local interaction matters, it asks when a \emph{publicly supported} device for organizing local interaction is socially justified in a thin regional ecosystem. In that sense, the paper is less about agglomeration in general than about the policy evaluation of a specific coordination instrument.

\subsection{Proximity, Lock-In, Regional Innovation Systems, and External Renewal}

A second relevant literature emphasizes that proximity has both benefits and risks. \citet{Boschma2005} and \citet{TorreRallet2005} argue that proximity can support learning and collaboration, but can also generate lock-in when interaction becomes overly repetitive. This concern is especially important for the present paper because its principal distortion is dynamic rather than static. The model formalizes the idea that dense current local interaction may reduce the novelty available to the regional system in the future.

The paper also speaks to the regional innovation systems literature, which stresses that regional performance depends not only on local scale, but also on institutional structure, organizational capacity, and system renewal \citep{CookeUrangaEtxebarria1997,AsheimLawtonSmithOughton2011,TodtlingTrippl2005,CoenenAsheimBuggeHerstad2017}. Work on path renewal and path creation similarly emphasizes that regions often need mechanisms that reconnect them to new knowledge combinations and outside opportunities \citep{ChaminadeBellandiPlecheroSantini2019}. In regional science terms, the key issue is not simply whether local interaction is dense, but whether a region can combine local coordination with access to external sources of renewal.

\citet{BatheltMalmbergMaskell2004}'s distinction between local buzz and global pipelines is especially useful here. The paper adopts that conceptual separation, but translates it into a policy comparison. The hub is modeled as a local coordination device inside the target region, whereas pipeline policy is modeled as a reduced-form mechanism that creates cross-regional renewal through outside links. This is the paper's main bridge to the knowledge-network and inter-regional collaboration literatures: it treats external linkages not as background context, but as a distinct policy margin that can be compared directly with local coordination in welfare terms. Empirical work on collaborative knowledge production supports the relevance of this distinction. \citet{HoekmanFrenkenVanOort2009} show that inter-regional knowledge collaboration in Europe is strongly shaped by geographic and institutional distance, while \citet{NeulaendtnerScherngellLiuJiDingPaier2025} find that the benefits of inter-city R\&D collaboration depend on place-based conditions and local knowledge endowments. Yet these literatures typically do not provide a welfare comparison among \emph{hub alone}, \emph{pipeline alone}, and the \emph{bundle}. That is the narrower gap this paper tries to fill.

The extension on participant composition also connects to work on absorptive capacity and anchor actors. Regions differ not only in headcount, but also in whether they contain actors that can transform contact into learning and collaboration \citep{CohenLevinthal1990,AgrawalCockburn2003}. That extension remains secondary to the baseline model, but it reinforces the regional-science point that the effective thickness of an ecosystem depends on composition as well as size.

\subsection{Place-Based Policy, Welfare Evaluation, and This Paper's Contribution}

A third relevant literature studies the welfare properties of geographically targeted policy. \citet{KlineMoretti2014ARE} provide a general welfare framework for local economic development policy, and \citet{KlineMoretti2014QJE} show in a major historical setting that place-based interventions can have durable local effects. More broadly, \citet{NeumarkSimpson2015} stress that the case for place-based policy depends on the underlying distortion and on the relevant margins of adjustment. That perspective is directly relevant here. A public hub is a place-based intervention, but its value depends on more than visible local activity. It depends on whether the intervention corrects a meaningful regional distortion and on how it interacts with external renewal.

The paper also speaks to a more specific literature on collaborative innovation policy. Here the relevant question is not only whether public support raises innovation, but whether policy-induced collaboration helps regions access external knowledge and diversify into new activities. The evidence is informative but not uniformly optimistic. \citet{BednarzBroekel2019} show that policy-induced R\&D networks need not automatically translate into stronger inter-regional knowledge diffusion, while \citet{MewesBroekel2020} find that subsidized joint R\&D projects can be more effective for regional technological diversification than subsidies to single organizations. For the present paper, the important point is that collaborative policy should be evaluated relative to the specific regional margin it changes. Existing work tells us that collaborative support can matter, but it does not solve the baseline problem studied here: whether a publicly supported hub should be judged against no policy or against a renewal policy that already builds outside links.

The paper contributes to this literature in three specific ways. First, it changes the relevant policy baseline: because pipeline policy has positive stand-alone value, the appropriate benchmark is often not hub versus no hub, but pipeline alone versus a bundle that combines local coordination with external renewal. Second, it introduces a clean regional distinction between \emph{local coordination} and \emph{cross-regional renewal} as separate policy margins. Third, it derives a welfare warning from a single principal wedge, a \emph{dynamic coordination externality}: visible local activity need not be an adequate metric when present interaction changes the future novelty of the regional system. The paper is therefore less a general theory of regional development than a stylized policy comparison targeted at one specific gap in the literature.

This contribution is deliberately stylized. The model does not attempt to explain endogenous network formation, migration, or location choice in full generality. Instead, it provides a tractable regional policy framework that isolates a single welfare wedge and shows how that wedge changes the evaluation of hub policy in thin regional ecosystems. That is the sense in which the paper is intended as a policy-relevant stylized framework for regional science rather than as a broad structural theory of regional development.
